\section{Complete tables, source interpretation and retained computation}
The companion \emph{The 137 surviving quotient classes} is an integral part of this common paper. Its two full tables contain every surviving historical label, subgroup order, both facet counts, all $349$ maximal-containment links, the actual computed subgroup descriptions and all normalizer quotients. It is included unchanged as \texttt{complete-census.pdf}; its five pages belong to the common acquisition, so neither application silently borrows another application's table. The exact subgroup generators and all eight forbidden-conjugate witnesses are supplied as data in addition to the readable tables. They are outputs of the exhaustive computation, not a supplied list on which the exhaustive algorithm depends.

\paragraph{What the source tables do and do not specify.}
The classification gives exact permutation subgroups, stronger data than semidirect-product abbreviations. The historical subgroup-description column is checked for order, while the full printed automorphism column is checked up to the elementary direct-product identities described in the census paper. We do not infer an abstract subgroup isomorphism merely from its abbreviated type and order. In particular source rows $38,39,40$ all use $19:9$-style notation, while the actual centers have orders $1,1,3$; the distinct action at row $40$ remains explicit. Exact representatives preserve all actions.

The accessible arXiv edition has these additional discrepancies, kept distinct from the erratum and from the unseen journal edition. Table~1's order-$70$ row labelled $58$ corresponds to row $63$ in the complete tables; their row $58$ instead has order $76$. Page~14 row $21$ prints $2^2\times(19:3)$, whereas its order is $456$ and its actual group is $2^3\times(19:3)$. Row $27$ prints $2\times19:3$, whereas its order is $342$ and its actual group is $2\times(19:9)$. Row $41$ has an elementary abelian normal subgroup of order eight and exponent two; its printed $8:(7:3)$ is not literally a cyclic-eight semidirect product under the source's stated notation. These exact finite distinctions, and the $38$--$40$ center distinction, are recorded separately. None is silently transferred to the journal text.

The source's numerical timing and former software history are preserved as historical reports, not mathematical results to reproduce. Its $145$, $30$, $69$, $70$ and $75$ census totals are replaced, respectively, by $137$, $26$, $65$, $66$ and $71$ in their stated categories. Its two named simple quotients, universal presentation, local types, covering scope and quotient invariants are all retained in the corrected principal contract.


\paragraph{Algorithms and their trust boundary.}
The complete common appendix contains the six distinct required files: the presentation calculation, marking construction, geometric tests, literal maximal-word evaluation, exhaustive subgroup census and Python table reconciliation. The shared marking file is printed once. The whole historical \texttt{source\_tables.json} is required comparison data. Neither the old unavailable \texttt{g.txt} nor a saved group workspace is used. GAP's exact group, subgroup, presentation and Tietze algorithms, ACE's completed coset enumeration, the named TomLib $J_1$ model and NetworkX's graph-isomorphism routine remain computational premises. Their actual versions and selected file hashes are retained; this does not claim a proof or a hash binding of the entire transitive installation. All complete code and data remain in the declared scope. There is no new Lean declaration or claim of end-to-end formal verification.

\paragraph{Attribution and proof organization.}
The principal universal group, geometry and source tables are due to Hartley and Leemans. The present edition reuses the already reviewed conjugacy tests, orbit counts, literal marking computation, finite-presentation upper bound and full corrected census. The new organization changes no finite calculation: the group and local geometry are established before the census, and both are now part of the shared proof. Thus the erratum application needs no unpriced theorem from the original-paper application. The two applications specify their distinct source contracts and draw their conclusions from this complete common paper, its table companion and its full required programs. The original fifteen-page preprint and two-page erratum remain distinct manuscripts.
